% 图契约：本题的源、边界和推导逐项对应；箭头符号按图注解释。
\documentclass[tikz,border=7pt]{standalone}
\usepackage{fontspec,xeCJK,amsmath,amssymb}
\IfFontExistsTF{Times New Roman}{\setmainfont{Times New Roman}}{\setmainfont{TeX Gyre Termes}}
\IfFontExistsTF{Heiti SC}{\setCJKmainfont[AutoFakeBold=1.5]{Heiti SC}}{\setCJKmainfont[AutoFakeBold=1.5]{FandolHei-Regular}}
\usetikzlibrary{arrows.meta,bending,calc,patterns,positioning}
\definecolor{ink}{HTML}{243745}
\definecolor{blue}{HTML}{275DAB}
\definecolor{green}{HTML}{177C63}
\definecolor{red}{HTML}{BD4935}
\definecolor{orange}{HTML}{AA6C16}
\definecolor{purple}{HTML}{715393}
\tikzset{every node/.style={font=\fontsize{9.3}{12.5}\selectfont,text=ink},
axis/.style={-{Stealth[length=4pt]},line width=.65pt,draw=ink!65},
vec/.style={-{Stealth[length=5pt]},line width=1pt,draw=blue},
short/.style={-{Stealth[length=3pt]},line width=.8pt},
guide/.style={dashed,draw=ink!40,line width=.5pt},
card/.style={rounded corners=2pt,draw=blue!65,fill=blue!4,align=center,text width=6.5cm,inner xsep=9pt,inner ysep=7pt},
title/.style={font=\fontsize{10.5}{14}\selectfont\bfseries,text width=15cm,align=center},
label/.style={font=\fontsize{8.4}{11}\selectfont,align=center},
note/.style={label,text width=6.2cm}}
\begin{document}
\begin{tikzpicture}[x=1cm,y=1cm]
\path[use as bounding box] (-8,-3.5) rectangle(8,3.5);
\node[title] at (0,3.15) {M2-063\quad 板间吸引使距离减小};
\begin{scope}[shift={(-4,0)}]
\draw[ink,line width=2pt] (-1.4,-1.5)--(-1.4,1.5) (1.4,-1.5)--(1.4,1.5);\node[above] at(-1.4,1.5){$+Q,V_1$};\node[above] at(1.4,1.5){$-Q,V_2$};\draw[<->,orange] (-1.4,-1.95)--(1.4,-1.95) node[midway,below] {$d$};\foreach\y in{-.7,0,.7}{\draw[vec] (-1.2,\y)--(1.2,\y);}
\draw[vec,red] (-1.4,1.1)--(-.5,1.1);\draw[vec,red] (1.4,1.1)--(.5,1.1);
\node[note] at (0,-2.85) {两个约束下 $F_d$ 均为负；固定电压时计入电源交换功。};
\end{scope}
\node[card,draw=blue!65,fill=blue!4] at (3.9,1.8) {$\displaystyle\frac{dC}{dd}=-\frac{\varepsilon_0A}{d^2}$};
\node[card,draw=green!65,fill=green!4] at (3.9,0) {$\displaystyle F_d|_Q=-\frac{Q^2}{2\varepsilon_0A}$};
\node[card,draw=orange!65,fill=orange!4] at (3.9,-1.8) {$\displaystyle F_d|_V=-\frac{\varepsilon_0AV^2}{2d^2}$};
\end{tikzpicture}\end{document}